% !TeX program = pdflatex
% !TeX encoding = UTF-8
\documentclass[a4paper,12pt]{article}
\usepackage[codigo,tikz]{relatorio}
\makeatletter\g@addto@macro{\UrlBreaks}{\UrlOrds}\makeatother  % permite quebrar URLs longas

% ── DADOS DO RELATÓRIO ───────────────────────────────────────────────────────
\disciplina{Sistemas Paralelos e Distribuídos}
\professor{Ricardo Rios}

\titulo{Logística Distribuída de Soros Antivenenos para Regiões com Conectividade Intermitente}
\cabecalho{Logística distribuída de soros antivenenos}

\autor{Cristiano Peniche Ceccon}
\autor{João Victor Mattos Lopes}
\autor{Thiago Cordeiro de Melo}

\data{\today}

\descricao{%
  Entrega 1 -- Proposta do projeto. Apresenta o problema, a motivação, o
  sistema distribuído que será desenvolvido, os resultados esperados e um
  protótipo inicial em Python que ilustra a ideia.}

\begin{document}

\iniciarPreTextual
\fazcapa
\fazfolhaderosto

\begin{resumo}
  O Brasil já descentralizou o soro antiveneno para unidades de saúde
  remotas, inclusive polos indígenas da Amazônia, e isso praticamente dobrou
  a proporção de pacientes tratados nas primeiras seis horas. A
  descentralização, porém, criou um problema novo: o estoque passou a estar
  espalhado por dezenas de pontos com internet intermitente, e ninguém tem
  uma visão confiável de onde há soro, de qual tipo, de qual lote e em que
  quantidade. Ao mesmo tempo, os casos notificados offline chegam com dias de
  atraso, o que atrasa a percepção de surtos e distorce a distribuição das
  cotas. Propomos um sistema distribuído \emph{local-first}, no qual cada
  posto mantém uma réplica do estado em tipos de dados replicados livres de
  conflito (CRDTs), sincroniza por anti-entropia quando há enlace e continua
  operando sem rede. Sobre esse estado replicado o sistema oferece notificação
  \emph{store-and-forward} de casos, detecção de aumento anormal de casos com
  alerta direcionado à região afetada, roteamento do paciente até o posto mais
  próximo que tem o soro adequado e uma análise paralela de criticidade dos
  postos. Um protótipo com oito nós comunicando-se por TCP mostra que as
  réplicas convergem após partições de vários dias sem perder notificações e
  expõe, de forma quantificável, o custo de decidir com informação
  desatualizada.
  \palavraschave{sistemas distribuídos; CRDT; tolerância a partição;
  soro antiveneno; vigilância epidemiológica; Amazônia}
\end{resumo}

\sumario

\iniciarTextual

% =============================================================================
\section{Proposta}
\label{sec:proposta}

\paragraph{Integrantes.} Cristiano Peniche Ceccon, João Victor Mattos Lopes e
Thiago Cordeiro de Melo.

\paragraph{O problema.} Acidentes por animais peçonhentos são de notificação
compulsória no Brasil; entre 2007 e 2019 foram quase 1,9 milhão de
notificações~\cite{completude2023}. A maioria dos casos é leve, mas entre os
graves a letalidade passa de 4\%~\cite{bjihs2023}, e o prognóstico piora
quando o soro é aplicado mais de seis horas depois do acidente. O soro existe,
é produzido no país e é distribuído gratuitamente pelo SUS: o problema é
\textbf{acesso oportuno}. Na Amazônia, pacientes chegam a usar vários meios de
transporte e viajar dias até uma unidade com soro~\cite{monteiro2020}. Por
isso o Ministério da Saúde vem descentralizando o soro para
\emph{Pontos Estratégicos de Antivenenos} (PEA), que eram cerca de 2.319 em
abril de 2026~\cite{nt127}, e o projeto SAVING levou o soro a 14 polos-base
indígenas no Amazonas~\cite{farias2025}.

Descentralizar, no entanto, transforma um estoque central em dezenas de
estoques pequenos, em locais com energia e internet intermitentes. Hoje não
há uma visão confiável e atualizada de \emph{onde há soro, de que tipo, de
qual lote e em que quantidade}; a reposição depende de notificações que,
justamente nas áreas mais isoladas, chegam atrasadas ou não chegam, o que
reduz a cota que essas áreas recebem. É um problema de \textbf{estado
replicado sob partição de rede frequente e prolongada}, consequência direta de
uma política pública já em execução.

\paragraph{Por que é interessante e desafiador.} Pelo teorema CAP, sob
partição é preciso escolher entre consistência e disponibilidade. Aqui a
escolha é clara: um posto isolado por uma semana não pode deixar de registrar
uma dispensação ou uma notificação porque o servidor está inacessível. Mas
disponibilidade sem coordenação tem custo: decisões tomadas com visão
desatualizada (remessas duplicadas, paciente encaminhado a um posto que já
zerou o estoque) e invariantes que não podem ser garantidas localmente
(estoque nunca negativo). Equilibrar esses efeitos num cenário real é o
desafio central do projeto.

\paragraph{O que vamos desenvolver.} Um sistema distribuído
\emph{local-first} com três tipos de nó (central estadual, polo regional e
posto), em que cada nó mantém uma réplica completa em CRDTs e sincroniza por
anti-entropia com os vizinhos sempre que há enlace. Sobre essa base:
\begin{enumerate}
  \item \textbf{inventário distribuído} por posto, tipo de soro e lote, com
        validade e dispensação FEFO (vence primeiro, sai primeiro);
  \item \textbf{notificação resiliente a desconexão}: o caso é registrado
        offline e entregue quando houver rede, sem perda nem duplicidade;
  \item \textbf{alerta de surto por região}: polos e central detectam aumento
        anormal de casos e o alerta é entregue aos nós da região afetada,
        disparando reforço de estoque;
  \item \textbf{roteamento multimodal}: dado o local do acidente, indicar o
        posto alcançável mais cedo que tem o soro necessário;
  \item \textbf{análise paralela de criticidade}: quais postos, se zerarem o
        estoque, deixam comunidades a mais de seis horas do tratamento.
\end{enumerate}

\paragraph{Resultados esperados.} (i) Um sistema funcional, demonstrável com
nós em processos ou contêineres separados e partições de rede reais;
(ii) uma avaliação experimental com métricas de convergência, atraso de
notificação, atraso de entrega do alerta, tráfego e quantidade de decisões
tomadas com informação desatualizada, sob diferentes padrões de
conectividade; (iii) medição do ganho de paralelismo na análise de
criticidade; (iv) um desenho reaproveitável para outros insumos
descentralizados com cadeia de frio, como vacinas. O sistema \emph{não} faz
diagnóstico nem decide conduta terapêutica: a conduta segue a síndrome
clínica e os protocolos do Ministério da Saúde.

% =============================================================================
\section{Contexto e motivação}
\label{sec:contexto}

\subsection{Dimensão do problema}

A Tabela~\ref{tab:sinan} resume as notificações no SINAN entre 2007 e 2019.
O escorpionismo é o agravo mais frequente no país, mas na região Norte o
quadro se inverte: o Norte concentrou 31,6\% dos acidentes ofídicos
nacionais, e no Amazonas cerca de 90\% deles são causados pela
\emph{Bothrops atrox}, a jararaca-da-Amazônia~\cite{diversitas2022}. Os
acidentes ofídicos na Amazônia brasileira custam pelo menos R\$~40 milhões
por ano aos cofres públicos~\cite{butantan40mi}.

\begin{table}[H]
  \centering
  \caption{Acidentes por animais peçonhentos notificados no SINAN, 2007--2019~\cite{completude2023}.}
  \label{tab:sinan}
  \begin{tabular}{lrr}
    \toprule
    \textbf{Tipo de acidente} & \textbf{Casos} & \textbf{Proporção} \\
    \midrule
    Escorpionismo        & 1.111.300 & 59,4\% \\
    Ofidismo (serpentes) &   386.938 & 20,7\% \\
    Araneísmo (aranhas)  &   373.224 & 19,9\% \\
    \midrule
    \textbf{Total}       & \textbf{1.871.462} & \textbf{100\%} \\
    \bottomrule
  \end{tabular}
\end{table}

\subsection{A descentralização já está em curso}

O projeto SAVING, da Fundação de Medicina Tropical Dr.~Heitor Vieira Dourado
(vinculada à UEA) com o Instituto Butantan e o Ministério da Saúde, levou o
soro a 14 polos-base em 7 Distritos Sanitários Especiais Indígenas do
Amazonas~\cite{farias2025}. Segundo o Instituto Butantan, a proporção de
pacientes tratados nas primeiras seis horas passou de 40,8\% para
75,8\%~\cite{butantan40mi}. O relato do programa mostra também as
restrições do ambiente: cadeia de frio que exige energia 24~horas com
gerador de reserva, transporte multimodal (avião até Tabatinga, depois barco)
e comunicação entre equipes de campo e especialistas feita principalmente
por WhatsApp, SMS e telefone~\cite{farias2025}.

Em 2026, a Nota Técnica Conjunta nº~127/2026 do Ministério da
Saúde~\cite{nt127} padronizou os PEA, com e sem armazenamento, e trouxe três
pontos que conversam diretamente com este projeto:
\begin{itemize}
  \item os PEA devem estar, como parâmetro, a até 1h30 da população-alvo;
  \item unidades com estoque devem monitorar validade, comunicar consumo,
        perdas e necessidade de reposição, e estar cadastradas no Sistema de
        Informação de Insumos Estratégicos (SIES);
  \item a relação de PEA, com o tipo de atendimento de cada um, deve ser mantida
        atualizada e enviada às centrais de regulação do SAMU e aos CIATox,
        e todos os profissionais devem saber onde ficam os PEA da região.
\end{itemize}
Ou seja: manter uma visão atualizada de onde há soro já é uma exigência
normativa. O que falta é um mecanismo que funcione onde a rede não funciona.

\subsection{O círculo vicioso da subnotificação}

O Ministério da Saúde distribui mensalmente as cotas de soro aos estados com
base em critérios epidemiológicos, isto é, nas notificações do
SINAN~\cite{nt127}. Há indícios de subnotificação na região Norte e alta
proporção de campos ignorados nas fichas~\cite{completude2023}. Onde o
acesso é pior, notifica-se menos; notificando menos, recebe-se menos soro. Um
mecanismo de notificação que funcione offline ataca a causa desse ciclo.

% =============================================================================
\section{Decisões de projeto}
\label{sec:decisoes}

\paragraph{Disponibilidade acima de consistência forte.} Bancos com
consistência forte (consenso, \emph{locks} distribuídos) exigem maioria
alcançável para aceitar escrita; um posto isolado ficaria inoperante.
Adotamos consistência eventual forte com CRDTs~\cite{shapiro2011}: toda
operação é aceita localmente, e o \emph{merge} das réplicas é comutativo,
associativo e idempotente, então réplicas que viram o mesmo conjunto de
atualizações convergem para o mesmo estado, em qualquer ordem de entrega.
É a filosofia \emph{local-first}~\cite{kleppmann2019}.

\paragraph{CRDTs escolhidos.}
\begin{itemize}
  \item \textbf{Saldo por (posto, lote)}: PN-Counter (entradas somam; dispensações,
        descartes e transferências subtraem). Os metadados do lote (tipo de
        soro, validade) são imutáveis.
  \item \textbf{Notificações de caso e alertas}: conjuntos que só crescem
        (G-Set), indexados por identificador único. Reentregas são inofensivas,
        o que dá \emph{store-and-forward} com semântica de entrega exatamente uma vez
        do ponto de vista da aplicação.
  \item \textbf{Alertas com identificador determinístico} (região + semana
        epidemiológica): se o polo e a central detectarem o mesmo surto, os dois
        alertas viram um só no merge.
\end{itemize}

\paragraph{Disseminação.} Anti-entropia \emph{push-pull} entre vizinhos da
topologia (algoritmos epidêmicos~\cite{demers1987}, como no
Dynamo~\cite{decandia2007}). Não há caminho fixo: a informação de um posto
isolado pode chegar à central por qualquer vizinho, inclusive, no futuro, por
um barco que passa pela comunidade levando um dispositivo (\emph{data
mule})~\cite{pentland2004,fall2003}.

\paragraph{O que o CRDT não resolve.} Convergência não implica invariantes:
dois dispositivos isolados no mesmo posto podem dispensar o mesmo frasco e o
saldo convergido fica negativo. O protótipo detecta esse caso; para
evitá-lo, avaliaremos contadores limitados (\emph{bounded counters}), que
repartem previamente o direito de decrementar entre réplicas~\cite{balegas2015}.

\paragraph{Dados mínimos.} A notificação replicada carrega apenas o necessário
para logística e vigilância (posto, região, data, agravo, gravidade, soro e
ampolas), sem dados pessoais, em linha com a LGPD. O registro clínico
continua nos sistemas oficiais (SINAN, e-SUS APS, SI-PNI, SIASI).

% =============================================================================
\section{Arquitetura proposta}
\label{sec:arquitetura}

A Figura~\ref{fig:arquitetura} mostra a topologia do cenário usado no
protótipo, inspirado na calha do Solimões. A hierarquia (central, polos,
postos) espelha a organização regionalizada do SUS e define quem abastece
quem, mas a sincronização é entre vizinhos, sem ponto único de falha: se a
central cair, polos e postos continuam trocando estado.

\begin{figure}[H]
  \centering
  \begin{tikzpicture}[
      font=\sffamily\footnotesize,
      no/.style={draw, rounded corners=2pt, minimum width=2.25cm, minimum height=0.8cm,
                 align=center, fill=#1},
      no/.default=white,
      bom/.style={thick},
      ruim/.style={thick, dashed},
      >=Stealth]
    \node[no=blue!12]  (man) at (0,0)      {Central\\Manaus};
    \node[no=green!12] (tab) at (-3.4,-1.9) {Polo\\Tabatinga};
    \node[no=green!12] (tef) at (3.9,-1.9)  {Polo\\Tefé};
    \node[no] (bc)  at (-5.6,-3.9) {Posto\\Benjamin Constant};
    \node[no] (ata) at (-3.4,-5.4) {Posto\\Atalaia do Norte};
    \node[no] (spo) at (-1.3,-3.9) {Posto\\S. Paulo de Olivença};
    \node[no] (alv) at (2.5,-3.9)  {Posto\\Alvarães};
    \node[no] (uar) at (5.3,-3.9)  {Posto\\Uarini};
    \draw[bom] (man) -- (tab);
    \draw[bom] (man) -- (tef);
    \draw[ruim] (tab) -- (tef);
    \draw[bom] (tab) -- (bc);
    \draw[ruim] (tab) -- (ata);
    \draw[ruim] (bc) -- (ata);
    \draw[ruim] (tab) -- (spo);
    \draw[bom] (tef) -- (alv);
    \draw[ruim] (tef) -- (uar);
    \draw[ruim] (alv) -- (uar);
    \node[draw, dashed, rounded corners, fit=(tab)(bc)(ata)(spo), inner sep=6pt,
          label={[font=\sffamily\scriptsize]below:Região Alto Solimões}] {};
    \node[draw, dashed, rounded corners, fit=(tef)(alv)(uar), inner sep=6pt,
          label={[font=\sffamily\scriptsize]below:Região Médio Solimões}] {};
    \begin{scope}[shift={(2.0,-5.6)}]
      \draw[bom] (0,0) -- (0.7,0) node[right] {enlace estável ($\geq$80\%)};
      \draw[ruim] (0,-0.45) -- (0.7,-0.45) node[right] {enlace intermitente};
    \end{scope}
  \end{tikzpicture}
  \caption{Topologia do cenário: cada linha é um par de nós que sincroniza
  diretamente quando o enlace está ativo. O alerta de surto do Alto Solimões
  é destinado apenas aos nós daquela região (e à central). Municípios reais,
  enlaces e tempos ilustrativos.}
  \label{fig:arquitetura}
\end{figure}

A Figura~\ref{fig:no} mostra os componentes de cada nó. Toda operação do
usuário (dispensar, receber remessa, notificar caso, consultar onde há soro)
é feita sobre a réplica local, portanto funciona offline. O detector de surto
roda nos polos (para a própria região) e na central (para todas).

\begin{figure}[H]
  \centering
  \begin{tikzpicture}[
      font=\sffamily\footnotesize,
      caixa/.style={draw, rounded corners=2pt, align=center, minimum height=0.85cm, fill=#1},
      caixa/.default=white,
      >=Stealth]
    \node[caixa=yellow!15, minimum width=3.6cm] (ui) at (0,1.9)
         {Interface local\\{\scriptsize dispensar, notificar, consultar}};
    \node[caixa=blue!10, minimum width=7.6cm, minimum height=1.25cm] (rep) at (0,0)
         {\textbf{Réplica local (CRDTs)}\\[2pt]
          {\scriptsize inventário PN-Counter por (posto, lote) \quad G-Set de casos \quad G-Set de alertas}};
    \node[caixa, minimum width=2.4cm] (ae) at (-2.75,-1.95) {Anti-entropia\\{\scriptsize push-pull TCP}};
    \node[caixa, minimum width=2.4cm] (sur) at (0,-1.95) {Detector de surto\\{\scriptsize C2 (polo/central)}};
    \node[caixa, minimum width=2.4cm] (rot) at (2.75,-1.95) {Roteador\\{\scriptsize grafo multimodal}};
    \node[caixa=gray!10, minimum width=1.8cm] (viz) at (-6.2,-1.95) {Vizinhos};
    \draw[<->, thick] (ui) -- (rep);
    \draw[<->, thick] (ae) -- (rep.south -| ae);
    \draw[<->, thick] (sur) -- (rep);
    \draw[->, thick] (rep.south -| rot) -- (rot);
    \draw[<->, thick, dashed] (viz) -- node[above, font=\scriptsize] {estado} (ae);
  \end{tikzpicture}
  \caption{Componentes de um nó. A réplica é a única fonte de verdade local;
  a anti-entropia a funde com a dos vizinhos quando há enlace.}
  \label{fig:no}
\end{figure}

\paragraph{Fluxos principais.}
\begin{enumerate}
  \item \textbf{Atendimento}: o posto procura o soro indicado (ou um substituto
        aceitável, como o antibotrópico-laquético para acidente botrópico) na
        réplica local e dá baixa por FEFO. Se não houver, consulta a réplica local
        e calcula a rota multimodal até o posto mais próximo com saldo suficiente.
  \item \textbf{Notificação}: o caso entra no G-Set local e se propaga por
        anti-entropia até a central, por qualquer caminho disponível.
  \item \textbf{Surto}: o polo conta os casos da região que já conhece e aplica o
        método C2 do EARS~\cite{hutwagner2003}, que compara o dia com uma linha de
        base móvel de 7 dias encerrada 3 dias antes. Como as notificações chegam
        atrasadas, os últimos dias são reavaliados a cada execução.
  \item \textbf{Reabastecimento}: a central abastece os polos, e cada polo os seus
        postos, com base na sua visão local; em alerta, o estoque-alvo da região
        dobra.
\end{enumerate}

% =============================================================================
\section{Trabalhos relacionados}
\label{sec:relacionados}

\paragraph{Sistemas oficiais.} O SIES acompanha a movimentação de insumos
estratégicos, inclusive soros, e é exigido dos PEA com
armazenamento~\cite{nt127}; o SINAN registra as notificações. São sistemas
de registro centralizados que dependem de conectividade (ou de digitação
posterior) para refletir o estado real. Nossa proposta não os substitui: é
uma camada operacional que funciona offline e pode alimentá-los quando houver
rede.

\paragraph{Sistemas logísticos offline-first.} O mSupply, sistema de
gestão logística de medicamentos muito usado em países de baixa e média
renda, tem uma versão \emph{offline-first} em que cada unidade opera sem
internet e sincroniza com um servidor central quando possível~\cite{msupply}.
O DHIS2 tem aplicativos de coleta que trabalham offline e sincronizam
depois~\cite{dhis2}. Ambos sincronizam em estrela com um servidor central e
resolvem conflitos por regras da aplicação; nosso foco é a troca entre
vizinhos (sem depender de um caminho até a central), o uso de CRDTs com
garantias formais de convergência e o acoplamento entre estoque, vigilância
de surtos e roteamento.

\paragraph{Acesso geográfico ao soro.} A Fiocruz estimou o tempo real de
acesso a serviços com soro cruzando a distribuição da população com a malha
de deslocamento por modal~\cite{fiocruz2024}. Esses estudos são análises
estáticas; o nosso roteamento usa o \emph{estoque atual} de cada posto, o que
muda a resposta quando um posto zera.

\paragraph{Fundamentos de SPD.} CRDTs~\cite{shapiro2011} e sua variante por
deltas~\cite{almeida2018}; software \emph{local-first}~\cite{kleppmann2019};
algoritmos epidêmicos~\cite{demers1987,decandia2007}; redes tolerantes a
atraso e \emph{data mules}~\cite{fall2003,pentland2004}; o teorema
CAP~\cite{gilbert2002}; e, na vigilância, o EARS~\cite{hutwagner2003} e a
estatística de varredura espacial~\cite{kulldorff1997}.

% =============================================================================
\section{Protótipo inicial}
\label{sec:prototipo}

Para dar concretude à proposta, implementamos um protótipo em Python, só com a
biblioteca padrão (pasta \texttt{prototipo/}). Ele sobe oito nós, cada um com
seu servidor TCP, e simula 30 dias no cenário da
Figura~\ref{fig:arquitetura}: acidentes (com perfil de agravos do Amazonas),
dispensações, remessas, enlaces que caem e voltam, uma queda de 9 dias em
Atalaia do Norte e um surto no Alto Solimões entre os dias 14 e 19. Os
valores do cenário (tempos, estoques, taxas) são ilustrativos.

\begin{lstlisting}
cd prototipo
python3 demo.py               # simulação completa (~15 s)
python3 -m unittest -v        # propriedades do merge, rotas e detector
python3 bench_paralelo.py     # speedup da análise de criticidade
\end{lstlisting}

\begin{table}[H]
  \centering
  \caption{Resultados de uma execução do protótipo (\texttt{saida\_exemplo.txt}).
  Execuções diferentes variam nos detalhes, pois o gossip depende do escalonamento real.}
  \label{tab:demo}
  \begin{tabularx}{\textwidth}{Xr}
    \toprule
    \textbf{Métrica} & \textbf{Valor} \\
    \midrule
    Casos simulados / notificações perdidas & 111 / 0 \\
    Atraso mediano da notificação até a central & 7,4 h \\
    Atraso máximo (Atalaia do Norte, isolada) & 175 h \\
    Alerta de surto: atraso até Benjamin Constant / S.~Paulo de Olivença & 5,4 h / 2,4 h \\
    Alerta de surto: atraso até Atalaia do Norte (isolada) & 141,5 h \\
    Tentativas de sincronização sem enlace & 42\% \\
    Tempo de convergência após restabelecer os enlaces & 35 ms \\
    Pacientes encaminhados a posto que já estava sem soro & 1 de 6 \\
    Estoque final de SAB em Atalaia (inicial: 40 ampolas) & 152 ampolas \\
    \bottomrule
  \end{tabularx}
\end{table}

A Tabela~\ref{tab:demo} mostra o comportamento esperado e o problema que
queremos atacar. Nenhuma notificação se perdeu e as réplicas convergiram
para o mesmo estado (mesmo \emph{hash}) logo que a rede voltou. Os testes
automatizados verificam, para sequências aleatórias de operações, que o merge
é comutativo, associativo e idempotente e que qualquer ordem de entrega
converge. Mas a mesma execução expõe o custo da visão desatualizada: sem
enxergar o consumo de Atalaia, o polo enviou remessas repetidas, e o posto
terminou com quase quatro vezes o estoque inicial, enquanto outro paciente
foi encaminhado a um polo que já não tinha o soro. Tornar visíveis as remessas
em trânsito e medir esses efeitos sob diferentes padrões de conectividade é
parte do trabalho das próximas entregas.

A Tabela~\ref{tab:bench} mostra a análise de criticidade N-1 num grafo
sintético com 22.500 localidades e 200 postos: para cada posto, um Dijkstra
com múltiplas origens sobre o grafo sem ele. As 200 tarefas são independentes
e foram distribuídas entre processos. O ganho satura em torno de 3,3 vezes
porque a máquina usada (Intel i7-12700H, notebook com 6 núcleos de desempenho
e 8 de eficiência) reduz a frequência sob carga: num laço puramente
aritmético, 16 processos rendem só 5,7 vezes. Caracterizar esse limite é um
dos pontos da avaliação.

\begin{table}[H]
  \centering
  \caption{Análise de criticidade N-1 em paralelo (22.500 vértices, 200 postos).}
  \label{tab:bench}
  \begin{tabular}{rrrr}
    \toprule
    \textbf{Processos} & \textbf{Tempo (s)} & \textbf{Speedup} & \textbf{Eficiência} \\
    \midrule
     1 & 13,00 & 1,00 & 100\% \\
     2 &  8,50 & 1,53 &  77\% \\
     4 &  5,00 & 2,60 &  65\% \\
     8 &  4,17 & 3,12 &  39\% \\
    16 &  3,90 & 3,33 &  21\% \\
    \bottomrule
  \end{tabular}
\end{table}

% =============================================================================
\section{Plano de trabalho}
\label{sec:plano}

\begin{table}[H]
  \centering
  \caption{Planejamento por entrega.}
  \label{tab:plano}
  \begin{tabularx}{\textwidth}{lX}
    \toprule
    \textbf{Entrega} & \textbf{Conteúdo} \\
    \midrule
    1. Proposta & Este documento e o protótipo de prova de conceito. \\
    2. Relatório v1 & Resumo, introdução, trabalhos relacionados e projeto;
      definição das métricas e dos cenários de conectividade. \\
    3. Relatório v2 & Um processo ou contêiner por nó, persistência local
      (SQLite), sincronização por deltas~\cite{almeida2018}, remessas em
      trânsito no estado replicado, contadores limitados; partições reais com
      \texttt{tc netem}; seção de implementação. \\
    4. Apresentação & Demonstração ao vivo: derrubar o enlace de um posto,
      registrar casos offline, disparar o alerta regional e mostrar a
      convergência ao religar. \\
    5. Relatório final & Avaliação experimental (convergência, atraso de
      notificação e de alerta, tráfego, decisões com visão desatualizada,
      speedup) e discussão. \\
    \bottomrule
  \end{tabularx}
\end{table}

% =============================================================================
\begin{thebibliography}{99}
\addcontentsline{toc}{section}{Referências}

\bibitem{completude2023} Completude das notificações dos acidentes por animais
  peçonhentos no SINAN: estudo descritivo, Brasil, 2007--2019.
  \textit{Epidemiologia e Serviços de Saúde}, 2023.
  \url{http://scielo.iec.gov.br/scielo.php?script=sci_arttext&pid=S1679-49742023000100303}.

\bibitem{bjihs2023} Análise epidemiológica dos acidentes com animais
  peçonhentos no Brasil, 2015--2020. \textit{Brazilian Journal of
  Implantology and Health Sciences}.
  \url{https://bjihs.emnuvens.com.br/bjihs/article/view/7820}.

\bibitem{diversitas2022} Epidemiologia de acidentes ofídicos no estado do
  Amazonas entre 2010--2020. \textit{Diversitas Journal}.
  \url{https://diversitasjournal.com.br/diversitas_journal/article/view/2138}.

\bibitem{nt127} BRASIL. Ministério da Saúde. \textit{Nota Técnica Conjunta
  nº 127/2026-SVSA/SAPS/SAES/SESAI/MS}: orientações aos gestores do SUS para
  a escolha dos Pontos Estratégicos de Antivenenos (PEA). Brasília, maio 2026.

\bibitem{farias2025} SEABRA DE FARIAS, A. et al. Decentralization of
  snakebite antivenom treatment to indigenous community health centers in the
  Brazilian Amazon: from demand to the first treatment (the SAVING Program).
  \textit{PLoS Neglected Tropical Diseases}, 2025.
  \url{https://pmc.ncbi.nlm.nih.gov/articles/PMC12043230/}.

\bibitem{monteiro2020} MONTEIRO, W. M. et al. Providing antivenom treatment
  access to all Brazilian Amazon indigenous areas: ``Every life has equal
  value''. \textit{Toxins}, 2020.
  \url{https://pmc.ncbi.nlm.nih.gov/articles/PMC7762137/}.

\bibitem{butantan40mi} INSTITUTO BUTANTAN. Envio de soro antiofídico para
  regiões remotas da Amazônia pouparia R\$ 40 milhões e 10 anos de vida de
  populações, diz pesquisa. \url{https://butantan.gov.br/noticias/envio-de-soro-antiofidico-para-regioes-remotas-da-amazonia-pouparia-r-40-milhoes-e-10-anos-de-vida-de-populacoes-diz-pesquisa}.

\bibitem{fiocruz2024} Estudo analisa acesso a soro antiofídico em regiões com
  alto índice de acidentes com cobras. Rádio Nacional / EBC.
  \url{https://radionacional.ebc.com.br/programas/tarde-nacional-amazonia/estudo-analisa-acesso-a-soro-antiofidico-em-regioes-com-alto-indice-de-acidentes-com-cobras}.

\bibitem{shapiro2011} SHAPIRO, M.; PREGUIÇA, N.; BAQUERO, C.; ZAWIRSKI, M.
  Conflict-free replicated data types. In: \textit{Symposium on
  Self-Stabilizing Systems (SSS)}, LNCS 6976, p.~386--400. Springer, 2011.

\bibitem{almeida2018} ALMEIDA, P. S.; SHOKER, A.; BAQUERO, C. Delta state
  replicated data types. \textit{Journal of Parallel and Distributed
  Computing}, v.~111, p.~162--173, 2018.

\bibitem{kleppmann2019} KLEPPMANN, M.; WIGGINS, A.; VAN HARDENBERG, P.;
  McGRANAGHAN, M. Local-first software: you own your data, in spite of the
  cloud. In: \textit{ACM Onward!}, 2019.

\bibitem{demers1987} DEMERS, A. et al. Epidemic algorithms for replicated
  database maintenance. In: \textit{ACM Symposium on Principles of Distributed
  Computing (PODC)}, 1987.

\bibitem{decandia2007} DeCANDIA, G. et al. Dynamo: Amazon's highly available
  key-value store. In: \textit{ACM Symposium on Operating Systems Principles
  (SOSP)}, 2007.

\bibitem{gilbert2002} GILBERT, S.; LYNCH, N. Brewer's conjecture and the
  feasibility of consistent, available, partition-tolerant web services.
  \textit{ACM SIGACT News}, v.~33, n.~2, 2002.

\bibitem{balegas2015} BALEGAS, V. et al. Extending eventually consistent
  cloud databases for enforcing numeric invariants. In: \textit{IEEE
  Symposium on Reliable Distributed Systems (SRDS)}, 2015.

\bibitem{fall2003} FALL, K. A delay-tolerant network architecture for
  challenged internets. In: \textit{ACM SIGCOMM}, 2003.

\bibitem{pentland2004} PENTLAND, A.; FLETCHER, R.; HASSON, A. DakNet:
  rethinking connectivity in developing nations. \textit{IEEE Computer},
  v.~37, n.~1, 2004.

\bibitem{hutwagner2003} HUTWAGNER, L. et al. The bioterrorism preparedness
  and response Early Aberration Reporting System (EARS). \textit{Journal of
  Urban Health}, v.~80, supl.~1, 2003.

\bibitem{kulldorff1997} KULLDORFF, M. A spatial scan statistic.
  \textit{Communications in Statistics -- Theory and Methods}, v.~26, n.~6,
  1997.

\bibitem{msupply} mSUPPLY FOUNDATION. \textit{Open mSupply: synchronisation
  (offline first)}. \url{https://msupply.foundation/open-msupply/sync}.

\bibitem{dhis2} DHIS2. \textit{DHIS2 Documentation}.
  \url{https://docs.dhis2.org}.

\end{thebibliography}

\end{document}
